\chapter{The general maxmin representation and its product}
\label{sec:p17-maxmin-representation}

This chapter gives the general finite-step representation of
Gilboa and Schmeidler, its uniqueness, the bounded-act extension,
and the product theorem \cite[pp.~143--152]{GilboaSchmeidler1989Maxmin}.
The state algebra need not be a sigma algebra, the outcome set
need not be finite, and priors are finitely additive.
The finite method of \cref{sec:p17-ambiguity-and-robust-choice}
is a special case, not a replacement for these domains.
We work in \(\mathsf{ZF}+\mathsf{AC}\), using the proved
Choice-to-Zorn implication \cref{thm:p06-ac-to-zorn}.

\section{The original domain and six axioms}

\begin{hbdefinition}{def:p17-mmeu-domain}{States, finite-step lottery acts and the source axioms}
Let \(S,X\ne\varnothing\), and let \(\Sigma\) be an algebra
of subsets of \(S\): it contains \(S\) and is closed under
complements and finite unions. Let \(Y=\Delta_f(X)\) be the
lotteries on \(X\) with finite support. Identify \(x\in X\)
with its certain lottery. The set \(L_0\) consists of acts
\(f:S\to Y\) constant on the cells of some finite
\(\Sigma\)-partition. Write \(y^*\) for the constant act
with value \(y\), and \(L_c=\{y^*:y\in Y\}\).
Mixtures of acts and lotteries are pointwise convex combinations.
For a preference \(\succeq\), use its strict part \(\succ\)
and indifference part \(\sim\); comparisons of lotteries mean
comparisons of their constant acts.

The source assumptions on \(L_0\) are:
\begin{enumerate}
\item[A.1] \(\succeq\) is complete and transitive.
\item[A.2] For \(f,g\in L_0\), \(h\in L_c\), \(0<t<1\),
\[
 f\succ g\quad\Longleftrightarrow\quad
          tf+(1-t)h\succ tg+(1-t)h .
\]
\item[A.3] If \(f\succ g\succ h\), there are \(a,b\in(0,1)\)
with \(af+(1-a)h\succ g\succ bf+(1-b)h\).
\item[A.4] If \(f(s)\succeq g(s)\) at every state, then
\(f\succeq g\).
\item[A.5] If \(f\sim g\), then \(tf+(1-t)g\succeq f\)
for every \(0<t<1\).
\item[A.6] Some two acts are strictly ranked.
\end{enumerate}
A.6 is the nondegenerate case; A.1--A.5 also permit universal
indifference. A.2 is stated in its original strict form.
Under A.1, applying it in both orientations also preserves
weak comparisons and indifference.
\end{hbdefinition}

\begin{hbtheorem}{thm:p17-mmeu-utility}{proved=annex:p17.mmeu.utility}{Lemma 3.1 on an arbitrary outcome set}
Under A.1--A.4 the constant-lottery preference has an affine
representation \(u:Y\to\mathbb R\), unique up to
\(u\mapsto au+b\), \(a>0\).
It has the form \(u(y)=\sum_x y(x)v(x)\) for a function
\(v:X\to\mathbb R\), which need not be bounded.
Under A.6 it is nonconstant, and can be normalized so that
its range \(K=u(Y)\) contains \([-1,1]\) in its interior.
This result uses the source's mixture continuity, not an
unstated topology on the outcome set.
\end{hbtheorem}

\begin{hbtheorem}{thm:p17-mmeu-certainty}{proved=annex:p17.mmeu.certainty}{Lemma 3.2 and exact certainty equivalents}
Fix such a \(u\). Every \(f\in L_0\) is indifferent to a
constant lottery. There is a unique number \(J(f)\in K\)
such that \(f\sim y^*\) implies \(J(f)=u(y)\).
The functional \(J\) represents \(\succeq\) and agrees with
\(u\) on constants. With strictly ranked upper and lower
constant bounds, the interpolating mixture coefficient is
unique; without strict bounds it is the utility number
that is uniquely determined.
\end{hbtheorem}

\section{The functional construction and its analytic premises}

\begin{hbdefinition}{def:p17-mmeu-function-spaces}{Simple functions, uniform completion and priors}
Let \(B_0\) be the real vector space of finite-valued
\(\Sigma\)-simple functions on \(S\), and let
\[
                     B=\overline{B_0}^{\|\cdot\|_\infty},
                \qquad \|a\|_\infty=\sup_{s\in S}|a(s)|.
\]
This uniform completion is used even when \(\Sigma\) is
only an algebra. Constants are denoted \(c\mathbf1\).
It is not obtained by replacing \(\Sigma\) by
\(\sigma(\Sigma)\).
Let \(\mathcal P(\Sigma)\) be the finitely additive
probabilities \(P:\Sigma\to[0,1]\), with \(P(S)=1\).
Give this set the topology of pointwise convergence on
events. A closed prior set means closed in this topology.
\end{hbdefinition}

\begin{hbtheorem}{thm:p17-mmeu-simple-functional}{proved=annex:p17.mmeu.simple.functional}{Lemma 3.3 with every extension specified}
Assume A.1--A.6 and normalize \(K\) as above.
There is a unique functional \(I:B_0\to\mathbb R\) obtained
from \(I(u\circ f)=J(f)\) by positive homogeneity.
It has
\[
\begin{gathered}
 I(c\mathbf1)=c,\qquad
 a\leq b\ \Longrightarrow\ I(a)\leq I(b),\\
 I(ta)=tI(a)\quad(t\geq0),\qquad
 I(a+b)\geq I(a)+I(b),\\
 I(a+c\mathbf1)=I(a)+c .
\end{gathered}
\]
In particular, for arbitrary \(a\in B_0\), choose \(M>0\)
with \(a/M\) valued in \([-1,1]\) and set
\(I(a)=M I(a/M)\). The answer is independent of \(M\).
Positive homogeneity does not assert oddness.
\end{hbtheorem}

\begin{hbtheorem}{thm:p17-mmeu-uniform-extension}{proved=annex:p17.mmeu.uniform.extension}{Lemma 3.4 and the completed function space}
The space \(B\) is complete. The functional satisfies
\[
                         |I(a)-I(b)|\leq\|a-b\|_\infty
\]
and has a unique continuous extension to \(B\), retaining
all displayed properties. For \(a_n\in B_0\) converging
uniformly to \(a\), it is \(I(a)=\lim_n I(a_n)\).
The space \(B\) is closed under products and under
continuous functions of any one bounded member.
\end{hbtheorem}

\begin{hbtheorem}{thm:p17-mmeu-hahn-banach}{proved=annex:p17.mmeu.hahn.banach}{The dominated linear extension used in separation}
Let \(E\) be a real vector space and \(p:E\to\mathbb R\)
be sublinear: \(p(x+y)\leq p(x)+p(y)\) and
\(p(tx)=tp(x)\) for \(t\geq0\).
Every linear functional on a subspace that is at most
\(p\) extends to a linear functional on \(E\) at most \(p\).
Consequently, for each \(b\in B\), there is a linear
functional \(L_b\) with
\[
                      L_b(a)\geq I(a)\quad(a\in B),
                               \qquad L_b(b)=I(b).
\]
Such \(L_b\) is positive, has \(L_b(\mathbf1)=1\),
and is continuous with norm one.
\end{hbtheorem}

\begin{hbtheorem}{thm:p17-mmeu-charges-compactness}{proved=annex:p17.mmeu.charges.compactness}{Integration and compactness without countable additivity}
Positive linear functionals \(L\) on \(B\) with
\(L(\mathbf1)=1\) correspond exactly to \(P\in\mathcal P(\Sigma)\):
\[
 P(E)=L(\mathbf1_E),\qquad
 L(a)=\int a\,dP .
\]
For simple \(a\), the integral is the finite partition sum;
for \(a\in B\), it is the limit of those sums under uniform
approximation. It is independent of the approximation and
bounded in absolute value by \(\|a\|_\infty\).

The space \(\mathcal P(\Sigma)\) is nonempty, compact and
Hausdorff. Each map \(P\mapsto\int a\,dP\), \(a\in B\),
is continuous. Thus on this normalized positive prior space
the eventwise topology is the topology generated by all
these integrals, the source's weak-star topology.
Every nonempty closed prior set attains the minimum of
every such integral. The proof constructs the compactness
premise explicitly from the ultrafilter extension principle.
\end{hbtheorem}

\begin{hbtheorem}{thm:p17-mmeu-support}{proved=annex:p17.mmeu.support}{Lemma 3.5 and the source's two separating sets}
For \(I(b)>0\), define
\[
 D_1=\{a\in B:I(a)>1\},\qquad
 D_2=\operatorname{conv}\bigl(\{a:a\leq\mathbf1\}
                          \cup\{a:a\leq b/I(b)\}\bigr).
\]
They are disjoint convex sets with nonempty interiors.
There is a positive normalized separator \(L_b\) such that
\[
               L_b(d_1)>1\geq L_b(d_2),\qquad L_b(b)=I(b).
\]
Choose its corresponding prior \(P_b\), and put
\[
       C=\overline{\operatorname{conv}}
                             \{P_b:I(b)>0\}.
\]
Then \(C\) is nonempty, closed and convex and, for every \(a\in B\),
\[
                              I(a)=\min_{P\in C}\int a\,dP .
\]
The same minimum representation holds for the intrinsic set
\[
              \{P\in\mathcal P(\Sigma):
                             \int a\,dP\geq I(a)\text{ for all }a\in B\}.
\]
\end{hbtheorem}

\begin{hblemma}{thm:p17-mmeu-finite-separation}{proved=annex:p17.mmeu.finite.separation}{The exact separation needed for prior uniqueness}
If \(C\subseteq\mathcal P(\Sigma)\) is nonempty, closed
and convex and \(P_0\notin C\), there is \(a\in B_0\) with
\[
                             \int a\,dP_0<\min_{P\in C}\int a\,dP .
\]
The function can be built from finitely many event indicators.
The proof uses a finite-coordinate neighborhood and a
proved Euclidean nearest-point separation, rather than
assuming countability of \(\Sigma\).
\end{hblemma}

\begin{hbtheorem}{thm:p17-mmeu-representation}{proved=annex:p17.mmeu.representation}{Theorem 1: both directions and exact uniqueness}
A relation on \(L_0\) satisfies A.1--A.5 if and only if
there are an affine \(u:Y\to\mathbb R\) and a nonempty
closed convex \(C\subseteq\mathcal P(\Sigma)\) such that
\[
 f\succeq g\quad\Longleftrightarrow\quad
       \min_{P\in C}\int u\circ f\,dP
                    \geq\min_{P\in C}\int u\circ g\,dP .
\]
The utility is unique up to positive affine change.
A.6 holds exactly when \(u\) is nonconstant; in that case
the closed convex set \(C\) is unique and equals the
intrinsic set in \cref{thm:p17-mmeu-support}.
If all acts are indifferent, \(u\) is constant and every
nonempty closed convex prior set represents the preference.
In this degenerate case \(C\) is unique only when
\(\Sigma=\{\varnothing,S\}\), for then there is only one prior.
\end{hbtheorem}

\section{The bounded-act extension on an algebra}

\begin{hbdefinition}{def:p17-mmeu-extension-domain}{The source domain and its convex uniform completion}
For a fixed nondegenerate constant-lottery order, let
\(L_{\rm lev}\) be the source's bounded level-measurable acts:
there are \(y_-,y_+\in Y\) with
\(y_+\succeq f(s)\succeq y_-\) for all \(s\), and both
\[
              \{s:f(s)\succ y\},\qquad\{s:f(s)\succeq y\}
\]
belong to \(\Sigma\) for every \(y\in Y\).
For a mere algebra this set need not be convex.
Define instead the containing convex domain
\[
 L_u=\{f:S\to Y:\ u\circ f\in B,\
               u(y_-)\leq u(f(s))\leq u(y_+)
                   \text{ for some }y_-,y_+\in Y\}.
\]
Both the source domain and this completion depend only on
the constant-lottery preference, not on its positive affine
utility normalization. In the degenerate case put \(L_u=Y^S\).
\end{hbdefinition}

\begin{hbtheorem}{thm:p17-mmeu-extension}{proved=annex:p17.mmeu.extension}{Proposition 4.1 with a valid convex domain}
One has \(L_0\subseteq L_{\rm lev}\subseteq L_u\), and
\(L_u\) is convex. The preference on \(L_0\) has a unique
extension to \(L_u\) satisfying A.1--A.5. In the
nondegenerate case it is
\[
                            f\succeq g\quad\Longleftrightarrow\quad
                                      I(u\circ f)\geq I(u\circ g).
\]
In the degenerate case every pair of extended acts is indifferent.
This gives its forced restriction to every original
level-measurable bounded act.
If \(\Sigma\) is a sigma algebra, \(L_{\rm lev}=L_u\).
For any convex act domain \(L\supseteq L_0\) carrying
A.1--A.5, the same formula is forced on \(L\cap L_u\).
Monotonicity alone is insufficient for this uniqueness.
\end{hbtheorem}

\begin{hbtheorem}{thm:p17-mmeu-domain-boundaries}{proved=annex:p17.mmeu.domain.boundaries}{The algebra, continuity and additive-measure boundaries}
The following examples give the stated failures.
\begin{enumerate}
\item Let \(S=\mathbb N_+\), \(\Sigma\) the finite/cofinite algebra,
and let utility be the high-prize probability in a binary lottery.
The acts with utility functions
\[
 a(n)=\tfrac12+\frac1{4n},\qquad
 b(n)=\tfrac12-\frac1{4n}+\frac{(-1)^n}{8n^2}
\]
belong to \(L_{\rm lev}\), but their half mixture does not:
its upper level set at \(1/2\) is the set of even integers.
Both acts and their mixture belong to \(L_u\).
\item A free ultrafilter on \(\mathbb N_+\) gives a finitely
additive prior \(P_{\mathcal U}(E)=\mathbf1_{\{E\in\mathcal U\}}\)
on all subsets. It has mass zero on every singleton and mass
one on the whole set, so is not countably additive.
Its integral of \(1/n\) is zero.
\item On the finite/cofinite algebra, ordering bounded utility
sequences by their ultrafilter order extends the simple-act
limit order and is complete and monotone, but ranks \(1/n\)
strictly above zero. It fails A.3, unlike the canonical
continuous extension.
\item On all subsets of \(\mathbb N_+\), finitely supported
priors form a convex nonclosed set. Their minimum on every
simple function equals that of all finitely additive priors.
Thus closedness cannot be omitted from the uniqueness claim.
\end{enumerate}
\end{hbtheorem}

\section{The geometric premises of the product theorem}

\begin{hbtheorem}{thm:p17-mmeu-extreme-geometry}{proved=annex:p17.mmeu.extreme.geometry}{Extreme points, projections and compact generating sets}
Let \(K\) be a nonempty closed convex subset of a prior space.
A face is a convex subset \(F\subseteq K\) such that an
interior convex combination in \(F\) of two members of \(K\)
forces both members into \(F\).
A point is extreme when its singleton is a face.
Then:
\begin{enumerate}
\item Every nonempty closed face has an extreme point of \(K\),
and \(K=\overline{\operatorname{conv}}\operatorname{Ext}(K)\).
\item An extreme point of a continuous affine image of \(K\)
in a Hausdorff real topological vector space
lifts to an extreme point of \(K\).
\item If \(D\) is compact and \(K=\overline{\operatorname{conv}}D\),
then \(\operatorname{Ext}(K)\subseteq D\).
\end{enumerate}
These results are proved for the eventwise topology, with
the needed compactness and separation supplied above.
\end{hbtheorem}

\begin{hbtheorem}{thm:p17-mmeu-exposed-points}{proved=annex:p17.mmeu.exposed.points}{The finite-dimensional exposed-point premise}
For a nonempty compact convex \(K\subseteq\mathbb R^d\),
\(K\) is the closed convex hull of its exposed points,
where an exposed point uniquely maximizes a linear functional.
For every linear functional \(v\cdot x\), exposed points
can approximate its maximum arbitrarily closely.
This is the precise finite-dimensional consequence required
in the source's appeal to Straszewicz.
\end{hbtheorem}

\begin{hbdefinition}{def:p17-mmeu-product}{Product algebras, product priors and the two independence clauses}
For nonempty \(S_i\), algebras \(\Sigma_i\), and nonempty
closed convex prior sets \(C_i\), put
\[
 S=S_1\times S_2,\qquad
 \Sigma=\Sigma_1\otimes\Sigma_2,
 \qquad
 C=\overline{\operatorname{conv}}\{P_1\otimes P_2:P_i\in C_i\}.
\]
Here \(\Sigma_1\otimes\Sigma_2\) is the algebra of finite
unions of rectangles, and \(P_1\otimes P_2\) assigns a
rectangle \(A\times B\) mass \(P_1(A)P_2(B)\).
A marginal act lifts by \(\bar f(s_1,s_2)=f(s_i)\).

For a represented nondegenerate preference with its unique
\(C\), two bounded utility acts \(f,g\) are independent in
the source sense when:
\begin{enumerate}
\item some \(P_0\in C\) minimizes both integrals
\(\int u\circ f\,dP\) and \(\int u\circ g\,dP\);
\item under every \(P\in\operatorname{Ext}(C)\), their
measurable event algebras are independent.
\end{enumerate}
Precisely, for a utility function \(a\), its measurable
event algebra consists of \(a^{-1}(D)\in\Sigma\) with
\(D\subseteq\mathbb R\) Borel. Clause 2 requires
\(P(A\cap B)=P(A)P(B)\) for every event from each of the
two such algebras. This makes the clause meaningful on
a mere algebra; it includes all the original strict and
weak level events and uses no undefined probability of
a set outside \(\Sigma\).
For sigma-algebra measurable functions this is the usual
full Borel-event independence clause.
\end{hbdefinition}

\begin{hbtheorem}{thm:p17-mmeu-product-priors}{proved=annex:p17.mmeu.product.priors}{The full finitely additive product construction}
Each product prior exists uniquely on the stated rectangle
algebra. The product map is continuous, its image is compact,
and \(C\) is nonempty, compact and convex.
The marginal images of \(C\) are exactly \(C_1,C_2\);
every extreme member of \(C\) is a product prior.
The two independence clauses are invariant under positive
affine utility changes. Every pair of lifted marginal acts
in \(L_u^1,L_u^2\) satisfies both clauses for this product.
\end{hbtheorem}

\begin{hbtheorem}{thm:p17-mmeu-product-representation}{proved=annex:p17.mmeu.product.representation}{Proposition 4.2: the complete product characterization}
Let two preferences on the marginal finite-step act domains
satisfy A.1--A.6 and agree on constant lotteries.
Use a common affine utility \(u\), their unique prior sets
\(C_1,C_2\), and their extensions to \(L_u^1,L_u^2\).
On the product completion \(L_u\), the preference represented
by the product \(C\) above is the unique preference satisfying:
\begin{enumerate}
\item A.1--A.6;
\item the marginal comparisons are preserved under their lifts;
\item every pair of lifted marginal acts satisfies both
source independence clauses.
\end{enumerate}
The conclusion holds for arbitrary state sets and algebras.
Its finite-subalgebra step is justified by extreme-point
lifts and does not replace the original state spaces.
\end{hbtheorem}

\begin{hbtheorem}{thm:p17-mmeu-product-examples}{proved=annex:p17.mmeu.product.examples}{Worked products and necessity of both independence clauses}
For each binary marginal, let the success probability range
over \([1/4,3/4]\). In state order \((00,01,10,11)\), the
product prior set is the convex hull of
\[
 (9,3,3,1)/16,\quad(3,9,1,3)/16,\quad
 (3,1,9,3)/16,\quad(1,3,3,9)/16 .
\]
Each is extreme and a product. Their convex hull also contains
\((5,3,3,5)/16\), which is not a product.
The marginal success bets each have value \(1/4\), and
their joint-success indicator has value \(1/16\).

Keeping only the middle two extreme products preserves the
marginal prior intervals and extreme-point stochastic independence,
but the two success bets have no common minimizing prior.
Conversely, with singleton fair marginals, the convex hull of
the diagonal and antidiagonal fair couplings has common minima
for all lifted marginal acts, but its extreme points are
not independent. Thus neither independence clause can be discarded.
\end{hbtheorem}
